In this section we present a simple proof for \autoref{thm:rank-competitiveness} by \cite{obmSimple}. Many proofs of this theorem use linear programming, but this one uses no special tools

\begin{lemma}
	\label{rank-cmp-compete-perfect}
	Competitiveness is determined by graphs with a perfect matching
\end{lemma}
\begin{proof}
	Let $G$ be a bipartite graph without a perfect matching, with orderings of offline agents $\sigma$ and of online agents $\pi$. Consider $G'$ which is a copy of $G$ without a vertex $x$ and orderings $\pi'$ induced by $V\backslash\{x\}$.

	Observe that $m=\mathrm{Ranking}(G,\pi,\sigma)$ being the matching found by the Ranking algorithm on $G$ with ranking $\pi$ and arrival order $\sigma$ is either identical to $m'=\mathrm{Ranking}(G',\pi',\sigma')$ or has one more vertex pair. This is apparent since either:
	\begin{enumerate}
		\item $x$ isn't in $m$,
		\item $x' = M(x)$ isn't matched in $m'$, or
		\item $x'$ is matched to $y$ in $m'$
	\end{enumerate}
	In the case of 3 there is an alternating path of edges in $m'$ and those that are in $m$ but not in $m'$. When this path ends in the same partition as $x$ then $m$ is larger than $m'$ by 1

	This also leads to the conclusion that if the respective matching sizes are different, they differ by  a single alternating path.

	Now consider the maximal matching of $G$ and the effect the removal of the vertices outside of it will have on the competitiveness.
	The size of the maximal matching will (of course) not change, while the Ranking's solution decreases by at most 1 with each removal, implying that this process of reducing $G$ to $G_p$ with a perfect matching didn't increase the competitiveness, meaning that analysis of graphs without perfect matchings is redundant.
\end{proof}

In the following section graphs are assumed to have perfect matchings. Let $M^\star$ be that perfect matching.

\begin{lemma}
	\label{rank-cmp-unmatched-simple}
	Let $v$ online vertex,$u=m^\star(v)$, $\sigma$ permutation.\\
	If $u$ isn't matched by $\mathrm{Ranking}(\sigma)$ then $v$ is matched by $\mathrm{Ranking}(\sigma)$ to $u'$ with a lower rank
\end{lemma}

\begin{proof}
	Trivial since $u,v$ are connected, therefore for $v$ to remain unmatched $u$ must've been matched to another vertex, which must've in turn had a lower rank than $v$
\end{proof}

\begin{lemma}
	\label{rank-cmp-unmatched-complicated}
	Let $v$ online vertex,$u=m^\star(v)$, $\sigma$ permutation, $\sigma_i$ permutation obtained by removing $v$ and inserting it into $sigma$ at the $i$-th place.\\
	If $u$ isn't matched by $\mathrm{Ranking}(\sigma)$ then for every $i$ $\mathrm{Ranking}(\sigma_i)$ matches $v$ with a vertex $u_i$ such that $\sigma_i(v_i)\leq\sigma(v)=t$
\end{lemma}

\begin{proof}
	$m=\mathrm{Ranking}(\sigma),m_i=\mathrm{Ranking}(\sigma_i)$.\\ Symmetrically to \autoref{rank-cmp-compete-perfect} $m,m_i$ are either identical or differ by a single alternating path. By repeatedly applying the logic from the proof of \autoref{rank-cmp-unmatched-simple} see that offline vertices in that path are of increasing rank in $\sigma_i$. Therefore $v_i=m_i(u),v'=m(u),\sigma_i(v_i)\leq\sigma_i(v')$\\
	Using \autoref{rank-cmp-unmatched-simple} we get $\sigma(v')<t$ and since $|\sigma_i(v')-\sigma(v')|\leq 1$ we receive $\sigma_i(v_i)<1+t$ equivalent to $\sigma_i(v_i)\leq t$ on integers
\end{proof}

The final lemma required to complete the proof uses this result directly.

\begin{lemma}
	\label{rank-cmp-probability}
	For random $\sigma$, the probability that the offline vertex of rank $t$ is matched by $\mathrm{Ranking}(\sigma)$ denoted $x_t$ fulfills the following: $1-x_t\leq\frac{1}{n}\sum\limits_{1\leq s\leq t}x_s$
\end{lemma}

\begin{proof}
	Let $\sigma$ be a random permutation of offline vertices. Define $\sigma'$ with the following random process: choose an offline vertex randomly with equal probability, take it out of $\sigma$ and insert it into the $t$-th rank. Let $u$ be the vertex of rank $t,u=m^*(u),R_{t-1}$\\ online vertices matched by the algorithm to offline vertices of rank $\leq t-1$\\
	probability that $v$ is not matched is $1-x_t$. Expcted cardinality of $R_{t-1}$ is $\sum\limits_{1\leq s\leq t-1}x_s$.\\
	Applying \autoref{rank-cmp-unmatched-complicated} with $\sigma:=\sigma',i$ chosen such that $\sigma'_i=\sigma$\ : if $v$ isn't matched by $\mathrm{Ranking}(\sigma')$ then $v$ is matched by $\mathrm{Ranking}(\sigma)$ to a vertex $u'$ such that $\sigma(v)\leq t$ implying $v\in R_t$\\
	Conditional on $\sigma$, $v\in R_t$ has a probability of $\frac{|R_t|}{n}$ since they are independent.\\
	We now have two random events $E,F$ that fulfill the condition that $E\implies F$. It is obvious to see that $\Pr[E]\geq\Pr[F]$\\
	The lemma follows
	
\end{proof}

With \autoref{rank-cmp-probability} the Theorem can be proved directly.\\
$G$ has a perfect matching implies that copetitiveness is $\inf \frac{1}{n} \sum\limits_{1\leq s\leq n}x_s$.\\
Let $S_t=\sum\limits_{1\leq s\leq t}x_s$. \autoref{rank-cmp-probability} can be rewritten as $S_t(1+\frac{1}{n})\geq 1 + S_{t-1}$ and the competitiveness ratio as $\frac{1}{n}S_n$.\\
Infimum occurs therefore when all inequalities are tight, resulting in $S_t=\sum\limits_{1\leq s\leq t}(1-\frac{1}{n+1})^s$ for all $t$, implying that competitiveness is at least $\frac{1}{n}\sum\limits_{1\leq s\leq n}(1-\frac{1}{n+1})^s=1-(1-\frac{1}{n+1})^n\xrightarrow[n\to\infty]{}1-\frac{1}{e}$

Thus the theorem is proven